\chapter{Conclusion and Future Work}
\label{ch:concl}

\section{Conclusion}

This thesis set out to measure, rather than assert, the workload relief that
an AI teaching assistant delivers, and to make such an assistant affordable
to operate at institutional scale. It presented a retrieval-augmented virtual
teaching assistant over NUS Canvas course materials, instrumented to estimate
the instructor-hours it returns, and integrated an extractive retrieval-
compression stage that, in an initial probe, removed roughly two-thirds of
prompt tokens on in-syllabus questions and correctly withheld context from an
out-of-syllabus query (\autoref{eval:sec:results:compression}).

\draftnote{final conclusion, once the instructor-hours and full
compression evaluations in \autoref{ch:evaluation} are complete.}

\section{Future Research Directions}

% TODO: fold the items below into the full future-work narrative once the
% conclusion is drafted.

\paragraph{Adaptive, evidentiality-guided compression.}
The extractive compression stage adopted from
\textsc{Recomp}~\cite{iclr24:Xu} carries two known limitations, identified
by~\citet{facl25:Jeong}~\cite{facl25:Jeong}: as the number of retrieved
documents grows, the compressor struggles to filter the increasing share of
irrelevant content, and a fixed compression ratio (here, the static threshold
$\tau$ and cap $k$ of \autoref{impl:sec:compression}) is suboptimal
per-question --- this system's own evaluation observed exactly this failure
mode when an equation-bearing sentence fell below the retention cut-off
(\autoref{eval:sec:results:compression}). Evidentiality-guided compression
with an adaptive per-query ratio, as in ECoRAG~\cite{facl25:Jeong}, is the
natural upgrade path, at the cost of an LLM-in-the-loop labelling pipeline
that the present system's hardware budget excluded
(\autoref{design:subsec:compression-choice}).

\paragraph{Broader evaluation of the compression stage.}
The token-reduction and latency results of
\autoref{eval:sec:results:compression} derive from a small probe over a
single course corpus; a larger course-material QA set across multiple
courses, with P50/P95 latency and answer-quality scoring
(\autoref{eval:subsec:tokenmetrics}), is required before the reductions can
be claimed as general.